\documentclass[12pt,letterpaper]{article}
\usepackage[margin=1in]{geometry}
\usepackage{times}
\usepackage{setspace}
\usepackage{footmisc}  % For footnotes

\begin{document}

\begin{center}
\textbf{\large SWORN DECLARATION OF ELVIS RYLAND NUNO}\\
\textbf{\large UNDER PENALTY OF PERJURY}\\
\vspace{0.3cm}
\textit{Submitted in Support of Criminal Referral}\\
\textit{State of Montana v. E'Lise Chard, et al.}
\end{center}

\vspace{0.5cm}

I, Elvis Ryland Nuno, declare as follows:

\section*{I. DECLARANT IDENTIFICATION AND BACKGROUND}

\begin{enumerate}

\item My name is Elvis Ryland Nuno. I am currently a resident of the State of Montana. I previously resided in Redmond, Washington from approximately December 2006 through December 2019, and in Missoula, Montana from approximately January 2017 through August 2020.

\item I am a senior telecommunications engineer with over 20 years of professional experience in telecom systems, ISP network architecture, and IT consulting. Prior to the events described herein, I maintained an unblemished professional reputation and earned between \$150,000--\$250,000 annually in contracts with major telecommunications companies and technology firms.

\item \textbf{Criminal History --- Complete Disclosure:} In the interest of transparency, I disclose all criminal matters of which I am personally aware and which I believe may appear in background checks. Given the inconsistent and often unreliable nature of criminal records systems—where records frequently vary depending on jurisdiction, vendor, and search methodology—this disclosure represents my good-faith effort to provide complete information based on my own knowledge and records:

\begin{itemize}
  \item \textbf{2007 --- Disturbing the Peace:} One minor misdemeanor ``disturbing the peace'' infraction when I was in my early twenties. This incident was unrelated to domestic violence, stalking, harassment, or any conduct alleged in subsequent cases involving Danielle Chard or E'Lise Chard.

  \item \textbf{March 2016 --- Edmonds Municipal Court --- DV Assault 3rd Degree (Coerced Plea):} I was compelled into accepting a guilty plea to a misdemeanor domestic violence assault in the third degree in Edmonds Municipal Court. This plea was not a reflection of my actual guilt, but the result of (1) multiple prosecutorial threats of more serious charges and potential incarceration; (2) ineffective assistance of counsel by attorney Patricia Fulton\footnote{See Exhibit 10, Section 2.12: WA\_State\_Bar\_Complaint\_Patricia-Fulton\_2016.pdf (formal bar complaint filed with Washington State Bar Association documenting incompetent/unethical representation).}, including failure to investigate exculpatory evidence and failure to challenge obvious inconsistencies in the complainant's statements; (3) the psychological impact and fear generated by the November 22, 2015 SWAT-style raid while sitting in my car at a gas station in Edmonds; and (4) my rapidly deteriorating financial ability to continue funding a private criminal defense after already spending tens of thousands of dollars.

  \item \textbf{September 2019 --- Washington Municipal Courts and Superior Court (Coerced Pleas):} In September 2019, I entered guilty pleas to misdemeanor no-contact order violations in Seattle Municipal Court, Lynnwood Municipal Court, and Edmonds Municipal Court, as well as misdemeanor stalking charges in Snohomish County Superior Court. These pleas were obtained through coercion arising from law enforcement threats against my romantic partner at the time (a widowed single mother) that she would ``lose custody of her child and never see her son again'' unless she cooperated and/or testified against me. This created an unconstitutional, coercive catch-22 described in detail in Section XI of this declaration.

  \item \textbf{Dismissed Cases --- Complete Vindication:} All other criminal charges filed against me in both Washington State and Montana have been dismissed with prejudice due to complete lack of evidence. This includes:

  \begin{itemize}
  \item \textbf{Seattle 2015--2016 Harassment and Cyberstalking:} Dismissed with prejudice in August 2016 after the Seattle Municipal Court found no evidence of criminal conduct and no credible basis for the charges brought by Danielle Chard.

  \item \textbf{Montana 2018--2019 Felony and Misdemeanor Stalking:} Both dismissed with prejudice after extended litigation and repeated failures by the State to produce any evidence supporting the alleged conduct. The Montana court explicitly recognized that my communications regarding E'Lise Chard---including my May 2018 YWCA complaint---constituted protected First Amendment speech and did not amount to criminal stalking or threats.
  \end{itemize}

\item \textbf{September 2016 --- Motion to Withdraw Coerced Plea:\footnote{See Exhibit 10, Section 1.3: Edmonds\_Trial-Nuno\_Declaration\_of\_Ineffective\_Assistance.pdf (formal declaration challenging plea voluntariness and asserting ineffective assistance of counsel by Patricia Fulton).}} In September 2016, recognizing that my March 2016 Edmonds Municipal Court plea had been involuntary and obtained under coercive conditions, I moved to withdraw my guilty plea. I based this motion on ineffective assistance of counsel, lack of a factual basis for the plea, and the coercive circumstances under which it was entered. The motion was denied.

\item \textbf{April 2017 --- Appeal of Withdrawal Denial:} With my last remaining financial resources I then retained appellate counsel and appealed the denial of my motion to withdraw the plea to Snohomish County Superior Court. In April 2017, the Superior Court denied the appeal. To the best of my understanding, this denial was based on the strict legal standard for withdrawing pleas post-sentencing, and not on any independent determination of my actual guilt.

\item \textbf{Significance of Withdrawal Attempts:} My persistent attempts to withdraw the 2016 plea---undertaken despite substantial financial strain, ongoing criminal accusations, and personal risk---demonstrate both my actual innocence and my belief that the plea was involuntary and constitutionally defective. These efforts are documented in the Edmonds Municipal Court and Snohomish County Superior Court records and are consistent with my position that I have been the target of a coordinated pattern of false accusations and prosecutorial overreach.

\item I submit this declaration in support of criminal investigation and prosecution of E'Lise Chard and, where applicable, Danielle Chard for perjury, criminal defamation, witness intimidation, and conspiracy to violate my civil rights. I am prepared to testify under oath before a grand jury or at trial regarding all matters contained herein, including the circumstances of the coerced pleas.

\end{itemize}
\end{enumerate}

\section*{II. RELATIONSHIP HISTORY WITH E'LISE CHARD}

\begin{enumerate}

\item I have known of E'Lise Chard since approximately 2002 through a few isolated social interactions starting with a very brief and hostile introduction in the early 2000s while I was dating her sister Danielle. Our first in-person meeting consisted of E'Lise stating to me: ``I hope you know I'm obligated to hate anyone my sister dates.'' We have never had a sustained in-person relationship, shared residence, or long-term in-person social contact.

\item Over the span of more than sixteen years, my in-person contact with E'Lise has amounted to only a handful of brief interactions, none of which involved threats, violence, or harassment on my part.

\item Our communications between approximately 2002 and 2017 were sporadic, primarily digital. There was never any regular or substantitive social interactions, platonic friendship, romantic relationship, and we never lived together.

\end{enumerate}

\section*{III. THE JUNE 2018 YWCA COMPLAINT (PROTECTED FIRST AMENDMENT SPEECH)}

\begin{enumerate}

\item On or about June 3, 2018, while residing in Missoula, Montana, I submitted a formal written complaint to YWCA Missoula President Deirdre Flaherty and Executive Director Cindy Weese regarding the conduct of YWCA employees E'Lise Chard and Rebecca Pettit.

\item My complaint was professional, detailed, and submitted in good faith. It alleged that E'Lise Chard and Rebecca Pettit had:

\begin{itemize}
\item Established a pattern of filing false police reports against me beginning in Washington State in 2016 and continuing in Montana through 2018;

\item Acted with vindictive intent stemming from E'Lise's resentment over my termination of a romantic relationship with her sister Danielle after Danielle became ``frighteningly abusive and violent'';

\item Contacted Missoula Police Officer Ethan Smith in June 2017 under the guise of official YWCA business;

\item Made knowingly false statements and presented fabricated evidence to Officer Smith;

\item Attempted to mischaracterize my professional IT career as evidence of involvement in a ``cyber crime spree'';

\item Abused their YWCA positions and organizational resources to create a facade of legitimacy for personal harassment.
\end{itemize}

\item My complaint detailed the consequences of their conduct, including Officer Smith's nearly year-long pattern of aggressive, off-the-record harassment of me and my parents, which contributed to severe mental health deterioration resulting in emergency psychiatric intervention for my mother.

\item I informed YWCA that my family had filed a formal complaint against Officer Ethan Smith and that I intended to pursue civil remedies if the harassment continued.

\item My complaint specifically requested that YWCA: (1) investigate the alleged misconduct; (2) take necessary protocol to ensure the activities ceased immediately; and (3) hold the responsible parties accountable for their behavior.

\item \textbf{As a gesture of good faith and support for YWCA's mission,} I enclosed a \$400 donation with my complaint letter, explicitly stating: ``I support the mission of your organization and appreciate the vital contribution to the community that it provides.''

\item My complaint was respectful, detailed, fact-specific, and requested institutional accountability---exactly the type of speech protected by the First Amendment right to petition for redress of grievances.

\item \textbf{YWCA never responded to my complaint.} They did not investigate, did not contact me, and took no action to address the alleged employee misconduct. Instead, they treated my letter and donation as suspected evidence of criminal actions and provided my letter and donation to E'Lise Chard to use when she filed a retaliatory protection order petition nine days later on June 12, 2018.

\item Montana courts subsequently found that my YWCA complaint constituted protected First Amendment speech and did not amount to criminal stalking, threats, or harassment. Judge Sam Warren (the presiding judge in protection order hearing) explicitly recognized the protected nature of this petition to a nonprofit organization regarding employee misconduct.

\end{enumerate}





\section*{III-A. THE 2025 MISJUSTICE ALLIANCE FORMAL NOTICES TO YWCA (CONTINUED INSTITUTIONAL INDIFFERENCE)}

\begin{enumerate}

\item On October 12, 2025, MISJustice Alliance (acting on behalf of myself and other victims of YWCA institutional corruption) sent a formal legal notice to YWCA Missoula requesting a formal response regarding systemic issues, possible criminal conduct, and conflicts of interest.\footnote{See Exhibit 10, Section 2.9: \texttt{Formal-Complaints/ywca-first-notice-10-12-2025.pdf}.}

\item The October 2025 notice specifically identified and requested explanation of:

\begin{itemize}
\item Detective Connie Brueckner's simultaneous role as a law enforcement officer and YWCA board member---an unprecedented conflict of interest that enabled the escalation of constitutional violations against me following my 2018 complaint;

\item Patterns of retaliation against complainants, including the documented retaliation I experienced;

\item Misuse of confidential client information, including HIPAA violations documented by other victims;

\item Funding practices that appear to prioritize crisis metrics over actual client recovery and stability;

\item The Board's conflict-management procedures, complaint-handling protocols, and response to allegations of structural institutional capture.
\end{itemize}

\item On December 5, 2025, MISJustice Alliance sent a comprehensive second notice to YWCA Missoula leadership, with copies to senior YWCA staff and key Missoula city and county officials, including the Mayor, Police Chief, and County Attorney.\footnote{See Exhibit 10, Section 2.10: \texttt{Formal-Complaints/ywca-second-notice-12-5-2025.pdf}.}

\item The December 2025 notice provided detailed documentation of:

\begin{itemize}
\item How E'Lise Chard used her YWCA position to import and amplify her sister Danielle's discredited accusations from Washington State;

\item How E'Lise recruited subordinate YWCA staffer Rebecca Pettit and YWCA liaison Officer Ethan Smith into the false narrative;

\item How this YWCA-coordinated campaign drove repeated, baseless arrests and harassment of myself and my family members;

\item How E'Lise allegedly recruited YWCA client Tyleen Root to conduct targeted online harassment and defamation campaigns against my business and customers;

\item How Missoula PD refused to intervene when I attempted to provide evidence of this harassment, and instead intimidated me---demonstrating a joint state-private conspiracy;

\item How YWCA ignored my 2018 formal complaint and a separate character-witness email about E'Lise's conduct;

\item How YWCA's institutional indifference destroyed my ability to continue sheltering vulnerable women I had been helping through my personal resources and network.
\end{itemize}

\item The December 2025 notice emphasized that this pattern of institutional conduct supports multiple federal causes of action, including:

\begin{itemize}
\item Civil rights violations under 42 U.S.C. \S{} 1983 (First, Fourth, and Fourteenth Amendment violations);

\item Civil rights conspiracy under 42 U.S.C. \S{} 1985 (coordinated deprivation of constitutional rights);

\item Defamation and intentional infliction of emotional distress under state tort law;

\item Civil and criminal RICO violations (18 U.S.C. \S{} 1962) based on pattern of racketeering activity;

\item Significantly increased civil and criminal exposure for YWCA Missoula and individual actors.
\end{itemize}

\item \textbf{To date---more than 60 days after the first notice and over one month after the second comprehensive notice---YWCA Missoula has provided absolutely no response.} No acknowledgment, no investigation, no dialogue, no corrective action.

\item This continued institutional silence in 2025, following detailed legal allegations copied to key governmental oversight officials, demonstrates the same pattern of deliberate indifference that enabled the decade-long harassment campaign against me:

\begin{itemize}
\item In 2018, YWCA ignored my formal complaint and instead provided my letter to E'Lise to use against me in a retaliatory protection order;

\item In 2025, YWCA ignores comprehensive legal notices documenting systematic constitutional violations and institutional corruption.
\end{itemize}

\item The absence of any response to formal legal notices establishes:

\begin{itemize}
\item Consciousness of institutional guilt;

\item Deliberate indifference to civil rights violations, supporting \textit{Monell v. Department of Social Services} liability;

\item Institutional policy of suppressing and ignoring complaints rather than addressing misconduct;

\item Continued obstruction of accountability and justice;

\item The complete futility of attempting internal remedies with a captured institution.
\end{itemize}

\item YWCA's refusal to respond demonstrates that the institution operates not as a victim-services nonprofit accountable to its charitable mission, but as a criminal enterprise that uses its nonprofit status and governmental partnerships to shield systematic abuse, suppress criticism, and enable coordinated harassment of individuals who dare to complain about employee misconduct or institutional corruption.

\end{enumerate}

\section*{IV. E'LISE CHARD'S JUNE 2018 PETITION, PERJURY, AND CONSTITUTIONAL VIOLATIONS}

\begin{enumerate}

\item On or about June 12, 2018---nine days after I filed my complaint with YWCA Missoula\footnote{See Exhibit 10, Section 2.8: Nuno-YWCA-Complaint.jpg (June 3, 2018 formal grievance letter with \$400 donation, requesting investigation of E'Lise Chard and Rebecca Pettit). Protection order filed June 12, 2018---nine days later, establishing temporal proximity for First Amendment retaliation.}---E'Lise Chard filed a petition for a temporary order of protection\footnote{See Exhibit 10, Section 1.4-1.5: ELise-OP-Filing-Page1.jpeg and ELise-OP-Filing-Page2.jpeg (sworn petition pages showing false statements contradicted by later testimony).} against me in Missoula County, Montana. In this petition, she alleged that I had engaged in a pattern of threatening, stalking, and harassing conduct toward her and her family, including supposed contact with her children.

\item In the written petition filed under penalty of perjury, E'Lise claimed that I had ``repeatedly'' contacted her, harassed her, and engaged in ongoing threatening behavior beyond the three emails I sent her in 2017. These written statements were sworn to be true.

\item At the subsequent protection order hearing, under oath and in open court, E'Lise was asked directly by the presiding judge whether I had made any contact with her beyond the three emails sent in 2017. She twice answered ``no'' under oath.

\item Her sworn testimony admitting that I had not contacted her beyond those three emails directly contradicted the written allegations in her petition that I had engaged in repeated, ongoing harassment and threats.

\item E'Lise's statements in the written petition and her sworn testimony cannot both be true. At least one version is knowingly false. This contradiction, together with additional inconsistencies documented in the court record, constitutes perjury under Montana law (MCA \S~45-7-202).

\item \textbf{I did not contest the protection order.} I stated to the court that I was fine with the order being issued, as I had no interest in or desire to have any contact with E'Lise Chard. I viewed the protection order as a simple civil matter that would formalize what was already my intent: to have no contact with her whatsoever.

\item I was unable to predict the unlawful legal assault that would follow.

\end{enumerate}

\section*{V. THE MONTANA STALKING PROSECUTION: EIGHTEEN MONTHS OF CONSTITUTIONAL VIOLATIONS}

\begin{enumerate}

\item Following the protection order hearing, criminal stalking charges were filed against me in Montana. These charges were based entirely on E'Lise Chard's false allegations and the compromised investigation by Officer Ethan Smith and Detective Connie Brueckner.

\item What followed was not a simple criminal case, but \textbf{eighteen months of severe, ongoing constitutional violations} that destroyed my career, health, and liberty prior to any conviction or finding of guilt:

\textbf{A. Search Warrant Based on Fabricated Evidence and Institutional Conflict}

\item Law enforcement executed a search warrant on my residence in Missoula, Montana, seizing electronic devices including computers, phones, tablets, and storage media essential to my telecommunications engineering work.

\item \textbf{The search warrant was obtained by Detective Connie Brueckner\footnote{\textbf{INSTITUTIONAL CONFLICT.} See Exhibit 10, Section 6.1: ywca-missoula-consolidated-financial-statements-years-2023-2024.pdf (ProPublica records confirm Detective Brueckner served on YWCA Missoula Board of Directors while investigating complaints against YWCA employees). This conflict was never disclosed to me, my defense counsel, or the court.}} and was based primarily on an affidavit incorporating investigative reports and statements from Officer Ethan Smith of the Missoula Police Department.

\item \textbf{Officer Ethan Smith's Evidence Was Fabricated:\footnote{See Exhibit 10, Section 2.11: Ty-Nuno-Police-Complaint-Ethan-Smith-03-2018.pdf (formal complaint by Elvis's father documenting Smith's year-long harassment pattern, fabricated reports, and mental health impact on family).}} As detailed in Section [cross-reference to Smith fabrications], Officer Smith's reports contained materially false statements, including:

\begin{itemize}
\item Fabricated phone conversations with my parents regarding my phone numbers and technical capabilities that my parents explicitly deny ever occurred;

\item False claims that my mother ``alluded to'' mental health issues, which both my parents categorically deny stating;

\item Mischaracterization of my professional IT skills as evidence of criminal intent rather than legitimate career expertise;

\item Reliance on E'Lise Chard's uncorroborated and contradictory allegations without independent verification.\footnote{See Exhibit 10, Section 2.2-2.3: MT-DOJ-POST-Complaint-Page1.jpeg and MT-DOJ-POST-Complaint-Page2.jpeg (formal POST complaint to Montana DOJ documenting fabricated evidence, conflict of interest, and misconduct by Smith and Brueckner).}
\end{itemize}

\item \textbf{Officer Smith was subsequently removed from my case} by Missoula Police Department due to ``obvious appearance of impropriety,'' specifically his close personal relationship with E'Lise Chard and questions about his objectivity in investigating complaints involving her. \textbf{This removal represented the sole instance of institutional accountability in the entire case, but was immediately followed by an unprecedented conflict of interest where a law enforcement officer with fiduciary duties to protect an organization was simultaneously investigating complaints against that organization's employees.}

\item \textbf{Detective Brueckner's Institutional Conflict Enabled Fraudulent Warrant:} Detective Connie Brueckner, who obtained the search warrant, simultaneously served on the Board of Directors of YWCA Missoula---the same organization where E'Lise Chard worked and where I had filed a formal complaint against E'Lise in June 2018.

\item Detective Brueckner was uniquely positioned to leverage her dual institutional roles to obtain the search warrant:

\begin{itemize}
\item \textbf{As YWCA Board Member:} She had insider knowledge of my June 2018 complaint against E'Lise, access to YWCA internal communications and records, and institutional loyalty to protect a YWCA employee from accountability.

\item \textbf{As Missoula Police Detective:} She had authority to seek search warrants, present evidence to judges, and coordinate with other officers (including the already-compromised Officer Smith).

\item \textbf{Leveraging Both Positions:} She could use the reputation and credibility of both YWCA (a respected nonprofit serving domestic violence victims) and Missoula Police Department (law enforcement authority) to present a compelling but fraudulent case to the issuing judge.

\item \textbf{Dual Institutional Weight:} A judge reviewing the warrant application would see: (1) a domestic violence services organization (YWCA) supporting the allegations; and (2) law enforcement (Missoula PD) vouching for the evidence---without knowing that the same person (Brueckner) sat on both sides and had an irreconcilable conflict of interest.
\end{itemize}

\item \textbf{The Search Warrant Was Based on Fraud:} Under \textit{Franks v. Delaware}, 438 U.S. 154 (1978), a search warrant is void if:

\begin{enumerate}
\item The affidavit contains deliberately false statements or statements made with reckless disregard for the truth; AND

\item The false statements were necessary to establish probable cause.
\end{enumerate}

Both elements are satisfied here:

\begin{itemize}
\item Officer Smith's fabricated statements (phone calls with my parents that never occurred, false mental health claims) were incorporated into the warrant affidavit;

\item These fabrications were necessary to establish probable cause, as there was no independent evidence supporting E'Lise's allegations;

\item Detective Brueckner knew or should have known of Smith's impropriety (later confirmed by his removal) and conflict of interest;

\item Detective Brueckner had her own undisclosed conflict (YWCA board membership) that violated her duty of candor to the court.
\end{itemize}

\item \textbf{Brueckner's Conflict Was Never Disclosed:} At no point during the warrant application process, the search execution, or the subsequent prosecution was Detective Brueckner's YWCA board membership disclosed to:

\begin{itemize}
\item The judge who issued the search warrant;
\item My defense counsel;
\item The prosecutors handling the case;
\item Me or my family.
\end{itemize}

This concealment violated Brady v. Maryland (duty to disclose exculpatory evidence) and constituted fraud upon the court.

\item \textbf{Seizure of Work Equipment Destroyed My Livelihood:} The search warrant authorized seizure of electronic devices, which included:

\begin{itemize}
\item Laptop computers containing proprietary client work, telecommunications network designs, and active project files;

\item Phones containing business contacts, client communications, and authentication credentials for work systems;

\item External storage devices with backup data, development environments, and years of professional work product;

\item Tablets and secondary devices used for client presentations and remote work.
\end{itemize}

\item The seizure of this equipment made it impossible for me to perform my job as a senior telecommunications engineer. I could not:

\begin{itemize}
\item Access client systems or networks;
\item Complete ongoing projects or deliverables;
\item Communicate with clients or employers using authenticated business accounts;
\item Maintain professional credentials requiring continuing education and active project work.
\end{itemize}

\item My employment was terminated as a direct result of the equipment seizure, and I was unable to secure new contracts or positions due to a flood of arrests and criminal charges that are returned from employment background checks.

\item \textbf{Equipment Was Held for Extended Period:} The seized equipment was retained by law enforcement for the entirety of 18 month criminal investigation, further compounding the economic harm and making it impossible to maintain my career.

\item When the equipment was eventually returned (after dismissal of all charges), much of the data and work product was obsolete, client relationships had been permanently severed, and my professional reputation had been irreparably damaged by the association with a criminal investigation---even though I was completely vindicated.

\item \textbf{The Fraudulent Warrant Directly Caused:}

\begin{itemize}
\item Loss of employment and income exceeding \$150,000 annually;
\item Destruction of client relationships and professional reputation;
\item Inability to complete contractual obligations, resulting in breach claims;
\item Loss of access to proprietary work product and professional portfolio;
\item Permanent exit from telecommunications engineering career;
\item Economic damages exceeding \$2.1 million over the 2018--2025 period.
\end{itemize}

\item \textbf{Legal Violations:} The fraudulent search warrant obtained by Detective Brueckner violates:

\begin{itemize}
\item Fourth Amendment (search based on false affidavit);
\item \textit{Franks v. Delaware} (knowing or reckless falsity in warrant affidavit);
\item Brady v. Maryland (concealment of Brueckner's conflict and exculpatory evidence);
\item Montana Rules of Professional Conduct for law enforcement (duty of candor to court, conflict of interest);
\item 18 U.S.C. \S~242 (deprivation of rights under color of law);
\item Montana criminal statutes regarding perjury and false statements to obtain warrants.
\end{itemize}

\end{enumerate}

\section*{VII. DETECTIVE CONNIE BRUECKNER: WEAPONIZING INSTITUTIONAL CONFLICTS}

\begin{enumerate}

\item Detective Connie Brueckner's simultaneous service on the YWCA Missoula Board of Directors while investigating criminal complaints against YWCA employee E'Lise Chard represents one of the most egregious conflicts of interest in this case.

\item \textbf{Brueckner's Dual Roles Created Unprecedented Leverage:}

\textbf{Timeline of Institutional Manipulation:}

\item \textbf{June 3, 2018:} I filed formal complaint with YWCA President identifying E'Lise Chard and Rebecca Pettit as sources of false police reports and harassment.

\item \textbf{June 12, 2018:} E'Lise filed retaliatory protection order petition (nine days after my YWCA complaint).

\item \textbf{June--September 2018:} Detective Brueckner, serving on YWCA board, assumed investigation after Officer Smith was removed for impropriety.

\item \textbf{As YWCA Board Member, Brueckner Had:}

\begin{itemize}
\item Access to my June 3, 2018 complaint against E'Lise;
\item Knowledge of YWCA's institutional interest in protecting employee E'Lise from accountability;
\item Awareness of the suppressed August 14, 2018 character reference letter from Witness PS (which YWCA never provided to investigators or my defense);
\item Institutional loyalty to YWCA's reputation and mission;
\item Fiduciary duty as board member to protect YWCA from liability;
\item Inside information about YWCA operations, employee matters, and organizational concerns.
\end{itemize}

\item \textbf{As Missoula Police Detective, Brueckner Had:}

\begin{itemize}
\item Authority to conduct criminal investigations;
\item Power to seek search warrants and present evidence to judges;
\item Credibility and presumption of law enforcement objectivity;
\item Access to police databases, investigative resources, and prosecutorial cooperation;
\item Ability to make arrest decisions and file criminal charges;
\item Professional relationships with prosecutors, judges, and other law enforcement.
\end{itemize}

\item \textbf{Brueckner Weaponized Both Positions:}

\begin{enumerate}
\item \textbf{Leveraged YWCA Reputation:} When seeking the search warrant and presenting evidence to prosecutors, Brueckner could implicitly or explicitly reference YWCA's involvement, lending credibility to E'Lise's allegations based on YWCA's respected role as a domestic violence services provider.

\item \textbf{Suppressed YWCA Evidence:} As a YWCA board member, Brueckner knew or should have known about the August 14, 2018 character reference letter from Witness PS that directly contradicted E'Lise's allegations. This exculpatory evidence was never disclosed.

\item \textbf{Protected YWCA from Accountability:} By aggressively pursuing charges against me (the person who filed a complaint against YWCA employees), Brueckner protected YWCA from institutional liability and validated YWCA's implicit decision not to investigate my June 3, 2018 complaint.

\item \textbf{Used Institutional Credibility to Obtain Fraudulent Warrant:} The judge reviewing Brueckner's warrant application would see a detective from a reputable police department, supported by allegations involving a respected domestic violence organization (YWCA), without knowing the detective sat on both sides.

\item \textbf{Prevented Internal YWCA Accountability:} My June 3, 2018 complaint requested YWCA investigate E'Lise's conduct. Brueckner's criminal investigation of me effectively shut down any possibility of internal YWCA accountability by framing me as the criminal rather than the complainant.
\end{enumerate}

\item \textbf{Brueckner's Conflict Should Have Required Recusal:}

Under basic law enforcement ethics and Montana POST standards, Detective Brueckner should have immediately recused herself from any investigation involving:

\begin{itemize}
\item YWCA employees (due to her board position);
\item Complaints filed against YWCA (I was a complainant);
\item Cases where YWCA had institutional interest in outcome;
\item Investigations where her dual roles created appearance of bias.
\end{itemize}

\item Instead, Brueckner not only failed to recuse, but actively concealed her YWCA board membership from all parties, including the court.

\item \textbf{This Concealment Was Deliberate:} Detective Brueckner's YWCA board membership was a matter of public record (documented in IRS Form 990 filings available on ProPublica). She knew of this conflict and deliberately failed to disclose it.

\item \textbf{Brueckner's Actions Constitute:}

\begin{itemize}
\item Fraud upon the court (concealing material conflict when seeking warrant);
\item Brady violation (suppressing exculpatory evidence);
\item Deprivation of civil rights under color of law (18 U.S.C. \S~242);
\item Conspiracy to violate civil rights (18 U.S.C. \S~241);
\item Professional misconduct (violation of Montana POST standards);
\item Abuse of authority (weaponizing institutional positions for personal/organizational interests).
\end{itemize}

\item \textbf{I Did Not Discover Brueckner's YWCA Conflict Until 2025:}

I only learned of Detective Brueckner's YWCA board membership in 2025 when researching this case for civil rights litigation. For over six years, this conflict was concealed from me, my defense counsel, prosecutors, and the courts.

\item This concealment prevented me from:

\begin{itemize}
\item Challenging the search warrant based on conflict of interest;
\item Moving to suppress evidence obtained through the fraudulent warrant;
\item Seeking Brueckner's recusal from the investigation;
\item Discovering and obtaining the suppressed Witness PS character reference;
\item Pursuing civil remedies while the underlying criminal case was pending;
\item Filing timely \S~1983 claims (statute of limitations arguments addressed separately).
\end{itemize}

\item The discovery of Brueckner's conflict in 2025 supports equitable tolling of statute of limitations for federal civil rights claims, as the conflict was fraudulently concealed and I could not have discovered it through reasonable diligence during the 2018--2019 criminal proceedings.

\end{enumerate}


\section*{VIII. CRIMINAL DEFAMATION BY E'LISE CHARD}

\begin{enumerate}

\item In her communications with law enforcement and in her June 2018 protection order petition, E'Lise Chard claimed that I had an ``extensive criminal drug history'' and that this was ``well known'' to police and judges.

\item This statement is categorically false. I have never been charged with, arrested for, or convicted of any drug-related offense in any jurisdiction.

\item E'Lise either knew this statement was false or acted with reckless disregard for the truth when she made it. Her statements were made to law enforcement, to the court, and to YWCA Missoula, and they have been used to destroy my reputation and career.

\item These false statements fit the elements of criminal defamation under Montana law (MCA \S~45-8-212): (1) a defamatory statement; (2) made knowingly or with reckless disregard for the truth; (3) communicated to third parties; (4) causing serious harm.

\end{enumerate}


\section*{IX. SUPPRESSION OF EXCULPATORY EVIDENCE BY YWCA AND DET. BRUECKNER}

\begin{enumerate}

\item In 2018, a mutual acquaintance (``Witness PS'') provided a written character reference\footnote{\textbf{CRITICAL EXCULPATORY EVIDENCE.} See Exhibit 10, Section 2.7: Nuno-PS-Witness-Statement-YWCA-Email-Complaint.pdf. This August 2018 testimonial describes my providing shelter to a pregnant homeless woman (Witness PS) at no cost in 2016, demonstrating my generous, non-violent character toward vulnerable women---the exact opposite of E'Lise Chard's false narrative and the exact mission YWCA claims to perform. YWCA and Detective Brueckner suppressed this evidence, never disclosing it to me, my counsel, or the court. This constitutes a Brady violation and institutional corruption.} on my behalf to YWCA Missoula, describing my non-violent nature and the implausibility of E'Lise's claims.

\item Detective Connie Brueckner of the Missoula Police Department, who conducted significant portions of the investigation into the Montana stalking allegations, was simultaneously serving on the board of YWCA Missoula.

\item This dual role created a profound institutional conflict of interest: a YWCA board member investigating criminal complaints about a YWCA-affiliated employee, where the complaint originated from a YWCA client (me) warning about that employee.

\item The suppression of Witness PS's exculpatory letter and the YWCA's failure to disclose this evidence constitute a violation of my due process rights under \textit{Brady v. Maryland}, as well as evidence of institutional corruption and conspiracy to protect E'Lise.

\end{enumerate}

\section*{X. COMPLETE JUDICIAL VINDICATION IN MONTANA AND WASHINGTON}

\begin{enumerate}

\item All felony and misdemeanor stalking charges brought against me in Montana were dismissed with prejudice after the State failed to produce evidence supporting E'Lise's allegations, despite having extensive access to my devices, communications, and digital records.

\item The 2015--2016 Seattle harassment and cyberstalking case was also dismissed with prejudice in August 2016, after Judge paid careful attention to the content and context of my online communications and correctly found that they were protected First Amendment speech, not criminal threats or stalking.

\item These dismissals with prejudice reflect a judicial determination that there was no evidence of criminal conduct and that the charges should not be refiled. They are consistent with my position that E'Lise and Danielle Chard repeatedly weaponized the legal system against me through false statements, with law enforcement and prosecutorial cooperation.

\end{enumerate}

\section*{XI. DAMAGES RESULTING FROM E'LISE AND DANIELLE CHARD'S CONDUCT}

\begin{enumerate}

\item As a direct and proximate result of the false accusations, criminal defamation, and perjury by E'Lise and Danielle Chard, I have suffered the near-total destruction of my professional career in telecommunications and network engineering.

\item I estimate that I have lost between \$6.4 and \$8.4 million in past and future earnings due to loss of contracts, inability to secure new employment, reputational damage, and forced relocation.

\item In addition to economic losses, I have suffered severe emotional distress, including symptoms consistent with post-traumatic stress disorder (PTSD), chronic anxiety, sleep disturbance, and social isolation.

\item These harms have been compounded by the length of the ordeal (2015--2025), the multi-jurisdictional nature of the coordinated attacks, and the ongoing fear of further malicious prosecution or retaliatory action.

\end{enumerate}

\section*{XII. PATTERN OF COORDINATION AND CONSPIRACY}

\begin{enumerate}

\item The actions of E'Lise and Danielle Chard did not occur in isolation. There is a clear pattern of coordination between them and multiple law enforcement officers and prosecutors across Washington and Montana, including but not limited to Edmonds Police, Seattle Police, Missoula Police, the Missoula County Attorney's Office, and municipal prosecutors in Edmonds, Lynnwood, and Seattle.

\item This pattern includes: (1) repeated filing of false and exaggerated reports; (2) reliance on those reports by law enforcement without proper verification; (3) suppression of exculpatory evidence; (4) retaliation when I attempted to expose their misconduct; and (5) the use of plea bargaining and threats to coerce admissions of guilt despite the lack of substantive evidence.

\item The coordinated nature of this conduct supports investigation under 18 U.S.C. \S 241 (conspiracy against rights), 18 U.S.C. \S 1512 (witness tampering), and parallel Montana criminal statutes.

\end{enumerate}

\section*{XIII. THE 2016 COERCED PLEA AND ITS AFTERMATH}

\begin{enumerate}

\item The March 2016 Edmonds Municipal Court plea occurred after the November 22, 2015 SWAT-style raid, which was triggered by Danielle Chard's false murder-suicide accusations and dramatically escalated the level of force and attention directed at me.

\item \textbf{The November 22, 2015 SWAT-Style Raid: Danielle's False Murder-Suicide Claims}

On November 22, 2015, while I was sitting in my car parked at a gas station in Edmonds, Washington, I was suddenly surrounded by a large number of police vehicles that rapidly converged on my location. Officers had their guns drawn and fixed on me, and commanded over a loudspeaker: ``ELVIS! TELL US WHERE YOUR GUN IS!''

\item This terrifying, life-threatening encounter was triggered by Danielle Chard's false report to law enforcement claiming that I was armed and intended to murder her and then commit suicide. These claims were completely fabricated:

\begin{itemize}
\item I did not have a gun---no weapon was found because no weapon existed
\item I had not made any threat of harm to Danielle or anyone else
\item I had not conducted myself in any way that would lead any reasonable person to infer I intended anyone harm
\item I posed no danger to Danielle, myself, or others
\item The entire ``armed and dangerous murder-suicide'' narrative was a malicious fiction created by Danielle Chard
\end{itemize}

\item \textbf{The Life-Threatening Risk Danielle Created}

The fact that I survived this encounter without being shot is itself remarkable. Police responses to reported armed suspects who have expressed intent to commit murder-suicide frequently result in officer-involved shootings. When officers believe they are confronting an armed, suicidal individual who has already stated intent to kill another person, the risk of lethal force is extraordinarily high.

\item Danielle's false report deliberately created a scenario where officers would perceive me as an immediate deadly threat, placing me at extreme risk of being shot. This was not merely a false accusation---it was a weaponization of law enforcement that could have easily resulted in my death.

\item No weapon was found. No evidence of threats was found. Yet this fabricated crisis---and the profound psychological trauma it inflicted---set the stage for the coerced plea that followed just four months later in March 2016

\item The combination of the SWAT raid, the criminal charges, and the financial burden of defending myself placed me under extreme psychological and economic pressure.

\item My retained counsel at the time failed to adequately investigate the allegations, challenge the evidence, or advise me of viable defenses, instead pressuring me to accept a plea to avoid the risk of a worse outcome at trial.

\item My subsequent attempts to withdraw that plea in 2016 and 2017 show that I did not and do not accept the legitimacy of that conviction and that I viewed it as coerced from the outset.

\end{enumerate}

\section*{XIV. THE 2019 WITNESS INTIMIDATION AND IMPOSSIBLE CATCH-22}

\begin{enumerate}

\item In 2019, multiple criminal cases were filed against me in Washington, including alleged no-contact order violations in Seattle Municipal Court, Lynnwood Municipal Court, and Edmonds Municipal Court, and alleged misdemeanor stalking in Snohomish County Superior Court. These charges were based largely on the complaints of Danielle Chard and related no-contact orders.

\item During this time, I was in a romantic relationship with a widowed single mother (``Partner A'') who had one young child. Partner A knew me well, was familiar with the history of false allegations by Danielle and E'Lise, and could have provided testimony favorable to my defense.

\item An officer with the Edmonds Police Department contacted Partner A and threatened her that if she did not testify against me or cooperate in the prosecution, she would ``lose custody of her child and never see her son again.''

\item This threat placed me in an impossible constitutional catch-22:

\begin{itemize}
  \item \textbf{Option A --- Expose the Threat:} If I disclosed this threat to the court and aggressively challenged the prosecutions, I reasonably believed law enforcement would follow through on their threat and use their leverage to initiate or influence custody proceedings against Partner A, thereby destroying her family and causing severe trauma to her child.

  \item \textbf{Option B --- Accept Coerced Pleas:} If I accepted guilty pleas, I could protect Partner A and her child from immediate retaliation and custody loss, but only by sacrificing my own constitutional rights and accepting criminal convictions based on fabricated or grossly exaggerated allegations.

  \item \textbf{No Constitutional Path:} There was no realistic third option that would both protect Partner A and her child and preserve my rights to a fair trial free from coercion.
\end{itemize}

\item Under these conditions, I accepted guilty pleas in September 2019 in the three municipal courts and in Snohomish County Superior Court. These pleas were not voluntary in any meaningful sense, but were the product of direct governmental coercion directed at a third-party witness and her child.

\item Law enforcement's threat against Partner A meets the elements of witness intimidation under 18 U.S.C. \S 1512 and analogous state statutes: (1) use of threats; (2) directed at a potential witness; (3) for the purpose of influencing testimony or cooperation; (4) under color of law.

\item I am prepared to identify Partner A to investigators and, with appropriate protections and assurances against retaliation or custody interference, she can corroborate the threats made by the Edmonds Police officer.

\item The 2019 witness intimidation is part of the broader pattern of conduct by E'Lise and Danielle Chard and collaborating officials. It demonstrates that when false accusations and weak evidence were not sufficient to obtain convictions, law enforcement escalated to coercive tactics that undermined the voluntariness of my pleas and violated my Fifth and Fourteenth Amendment rights.

\end{enumerate}


\section*{XVII. SLAPP PRINCIPLES: WEAPONIZATION OF LEGAL PROCESS ACROSS JURISDICTIONS}

\begin{enumerate}

\item Strategic Lawsuits Against Public Participation (SLAPP) refer to legal actions filed not to win on the merits, but to burden defendants with litigation costs and intimidate them into silence. While traditionally associated with civil defamation suits, the underlying principle---using legal process as a weapon to chill protected speech---applies equally to retaliatory criminal prosecutions.

\item The coordinated campaign against me across Montana and multiple King County, Washington jurisdictions exhibits every characteristic of a criminal SLAPP prosecution designed to punish my protected First Amendment speech.

\textbf{A. Classic SLAPP Characteristics Present in My Case}

\item \textbf{Immediate Timing Following Protected Speech}

On June 3, 2018, I filed a formal written complaint with YWCA Missoula documenting false police reports and harassment by E'Lise Chard and Rebecca Pettit. This complaint constituted core First Amendment petition activity.

\item Nine days later, on June 12, 2018, E'Lise filed a retaliatory protection order petition containing perjurious statements. Following Judge Warren's finding that my YWCA complaint was protected speech, criminal stalking charges were immediately filed---classic SLAPP timing designed to punish constitutionally protected activity.

\item \textbf{Multi-Jurisdictional Forum Shopping and Charge Stacking}

The harassment campaign deliberately weaponized legal process across multiple jurisdictions simultaneously:

\begin{table}[h]
\begin{tabular}{|l|l|l|}
\hline
\textbf{Jurisdiction} & \textbf{Charges Filed} & \textbf{Strategic Purpose} \\
\hline
Missoula Municipal Court & Multiple stalking charges & Local pressure, GPS monitoring \\
Missoula Superior Court & Felony stalking charges & Escalation, increased bail \\
Seattle Municipal Court & No-contact violations & Multi-state coordination \\
Lynnwood Municipal Court & No-contact violations & Resource drain \\
Edmonds Municipal Court & No-contact violations & Overwhelm capacity \\
Edmonds Superior Court & Misdemeanor stalking & Prosecutor coordination \\
Edmonds Superior Court & Felony stalking & Prosecutor coordination \\
\hline
\end{tabular}
\caption{Coordinated Multi-Jurisdictional SLAPP Prosecution}
\end{table}

\item This deliberate forum shopping forced me to retain five separate attorneys across two states and seven different courts, creating an insurmountable financial burden designed to prevent effective defense and punish protected speech.

\item \textbf{Dismissal and Re-Filing to Escalate Bail}

Prosecutors in Missoula engaged in the classic SLAPP tactic of dismissing charges and immediately re-filing them with elevated charges to justify dramatically increased bail amounts:

\begin{itemize}
\item Initial arrest: \$10,000 bail (one misdemeanor stalking charge)
\item Second arrest: \$50,000 bail (one misdemeanor, two felony stalking charges)
\item Third attempt: Motion to increase bail to \$100,000 (ultimately unsuccessful)
\end{itemize}

\item This pattern demonstrates the use of criminal process not to pursue justice, but to maximize pretrial punishment and financial devastation---a hallmark of SLAPP prosecutions.

\item \textbf{Resource Drain and Economic Destruction}

The coordinated SLAPP campaign achieved its intended purpose of economic annihilation:

\begin{itemize}
\item Loss of \$150,000 annual IT consulting contract (June 2018)
\item Legal fees exceeding \$387,000 across multiple jurisdictions
\item Seizure of all work equipment during home raid, preventing contract completion
\item GPS monitoring restrictions preventing travel to fulfill work obligations
\item Loss of \$60,000 annual St. Regis School District contract (2024)
\item Total career destruction: \$2.1 million in lost earnings (2018--2025)
\end{itemize}

\item \textbf{Complete Lack of Merit: All Charges Dismissed}

After 18 months of Montana prosecution involving extensive investigation, seizure of all electronic devices, and analysis of communications:

\begin{itemize}
\item All Montana felony stalking charges: Dismissed with prejudice
\item All Montana misdemeanor stalking charges: Dismissed with prejudice
\item All Montana criminal defamation charges: Dismissed with prejudice
\item Seattle 2015--2016 harassment/cyberstalking: Dismissed with prejudice
\end{itemize}

\item The universal dismissal of charges with prejudice---meaning they cannot be refiled---demonstrates that no probable cause ever existed. These prosecutions were SLAPP actions from inception, designed solely to punish protected speech and silence institutional criticism.

\item \textbf{Chilling Effect on Protected Speech}

The SLAPP prosecution achieved its intended chilling effect:

\begin{itemize}
\item I was forced to flee Washington \& Montana to escape ongoing harassment
\item I have been effectively silenced from making further complaints about institutional misconduct
\item My family developed severe agoraphobia and PTSD from fear of law enforcement
\item Future complainants are deterred from reporting YWCA misconduct
\item The YWCA has never investigated my complaint or implemented any reforms
\end{itemize}

\textbf{B. Montana's 2025 Anti-SLAPP Law and Retrospective Application}

\item Montana's House Bill 292, enacted May 1, 2025, adopts the Uniform Public Expression Protection Act (UPEPA). While UPEPA cannot retroactively dismiss past prosecutions, it provides powerful evidence for federal civil rights claims by establishing that Montana now formally recognizes SLAPP prosecutions as violations of fundamental rights.

\item The SLAPP framework transforms my case from isolated incidents into a pattern of weaponizing legal process to silence protected speech---precisely what anti-SLAPP laws exist to prevent and punish.

\textbf{C. SLAPP as Evidence of Constitutional Violations}

\item The SLAPP characteristics documented here support multiple federal civil rights theories:

\begin{itemize}
\item \textbf{42 U.S.C. § 1983 First Amendment Retaliation}: YWCA complaint constituted core political speech; immediate criminal charges establish causal connection under \textit{Hartman v. Moore}
\item \textbf{Malicious Prosecution}: Universal dismissals prove lack of probable cause; SLAPP characteristics demonstrate malice
\item \textbf{Abuse of Process}: Dismissal/re-filing to increase bail demonstrates use of process as punishment rather than pursuit of justice
\item \textbf{Conspiracy (42 U.S.C. § 1985)}: Coordination between YWCA employees, police, and prosecutors across state lines establishes conspiracy to violate civil rights
\end{itemize}

\end{enumerate}

\section*{XVIII. WARRANTLESS ARREST AND UNLAWFUL INTERSTATE COORDINATION}

\begin{enumerate}

\item In August 2018, while I was on GPS monitoring and restricted to Missoula County under court-ordered release conditions, I was subjected to a warrantless arrest based on fabricated claims of a Washington State extradition warrant. This incident demonstrates the dangerous back-channel coordination between law enforcement agencies across state lines, weaponized at the request of private citizens to continue my harassment and incarceration.

\textbf{A. Court-Ordered Release Conditions and GPS Monitoring}

\item Following my attorney Bryan Tipp's successful negotiation of my release from Missoula County Jail, I was placed under the following restrictive conditions:

\begin{itemize}
\item GPS ankle monitoring with continuous location tracking
\item Movement restricted to Missoula County only---could not travel outside county without prior court approval
\item Exclusion zone: prohibited from coming within 1,500 yards of the YWCA office
\item This exclusion zone significantly restricted my movement within the city of Missoula itself
\end{itemize}

\item These restrictions made it impossible for me to physically appear in Washington State courts for arraignment on pending charges. My presence in Washington required extensive coordination between:

\begin{itemize}
\item Bryan Tipp (my retained Montana attorney)
\item Tim Milios (my retained Washington attorney)
\item Multiple Washington Public Defenders offices (for jurisdictions where I could not afford private representation)
\item Missoula County Pretrial Services
\item Various Washington courts (Seattle, Lynnwood, Edmonds Municipal, Edmonds Superior)
\end{itemize}

\item Bryan Tipp had already notified all Washington courts of my GPS monitoring restrictions and had scheduled coordinated appearance dates that would allow me to appear with prior court approval and travel authorization.

\textbf{B. E'Lise's Communication to Danielle and Fabricated Extradition Claim}

\item After learning that I had been able to post bail and was on GPS monitoring (information E'Lise obtained through her close coordination with Missoula law enforcement and YWCA staff appearing at my arraignment hearings), E'Lise immediately communicated this information to her sister Danielle Chard in Washington State.

\item According to Danielle, she would ``never be safe until [I] was behind bars.'' Acting on this stated goal, Danielle contacted an Edmonds police officer with whom she had an established relationship from prior false reports.

\item This Edmonds officer then made an \textbf{unofficial, back-channel call} to a Montana law enforcement officer, falsely claiming that Washington State intended to extradite me and that there was an outstanding warrant for my arrest.

\textbf{C. The Warrantless Home Entry and Arrest}

\item Based solely on this unofficial communication---with no actual warrant, no extradition request, and no formal interstate process---Montana officers entered my home without a warrant and arrested me.

\item The purported basis for the arrest was that I had failed to appear in Washington courts and that Washington prosecutors had called for my arrest and extradition.

\item \textbf{This was categorically false:}

\begin{itemize}
\item No Washington prosecutor had called for my arrest
\item No Washington prosecutor had requested extradition
\item No new warrants had been issued in any jurisdiction
\item My failure to appear was due to GPS monitoring restrictions, which had been communicated to all courts
\item Scheduled appearance dates had already been coordinated through my attorneys
\end{itemize}

\textbf{D. ``Lost in Jail'' for Multiple Days}

\item As a direct result of this unlawful arrest based on fabricated information, I was ``lost in jail'' for several days.

\item Jail officials could not release me because they did not have documentation of why I was in custody or what charges I was being held on.

\item There was no warrant, no extradition request, no new charges---nothing in the system to explain my detention.

\item In order to secure my release, my attorneys had to provide the court with proof that:

\begin{itemize}
\item No prosecutor in Montana had called for my arrest
\item No prosecutor in Washington had called for my arrest or extradition
\item No new warrants had been issued in any jurisdiction
\item My detention was unlawful and based on fabricated claims
\end{itemize}

\textbf{E. Constitutional Violations}

\item This warrantless arrest violated multiple constitutional protections:

\begin{itemize}
\item \textbf{Fourth Amendment}: Warrantless home entry and arrest without probable cause
\item \textbf{Fourth Amendment}: Arrest based on unverified, unofficial communication rather than lawful warrant
\item \textbf{Fourteenth Amendment Due Process}: Detention without charges or lawful basis
\item \textbf{Interstate Compact violations}: Circumvention of formal extradition process required by U.S. Constitution
\end{itemize}

\textbf{F. Evidence of Coordinated Interstate Conspiracy}

\item This incident demonstrates direct coordination and weaponization of law enforcement across state lines through back-channeled, unofficial communications:

\begin{table}[h]
\begin{tabular}{|l|l|}
\hline
\textbf{Participant} & \textbf{Role in Conspiracy} \\
\hline
E'Lise Chard (Missoula) & Monitored my custody status through MPD contacts \\
& Communicated my release to sister Danielle \\
\hline
Danielle Chard (WA) & Contacted Edmonds officer with false urgency \\
& Stated goal: keep Elvis ``behind bars'' \\
\hline
Edmonds Officer (WA) & Made unofficial call to Montana officer \\
& Manufactured false extradition claim \\
\hline
Montana Officer & Conducted warrantless home entry \\
& Arrested without warrant or probable cause \\
\hline
Missoula County Jail & Held without charges or documentation \\
& Could not explain basis for custody \\
\hline
\end{tabular}
\caption{Interstate Back-Channel Coordination for Unlawful Arrest}
\end{table}

\item This coordination bypassed all formal legal process and constitutional protections, demonstrating a pattern of law enforcement officers acting at the behest of private citizens to achieve personal vendettas rather than legitimate law enforcement purposes.

\textbf{G. Pattern of Abuse: Back-Channel Communications}

\item This warrantless arrest is part of a broader pattern of unofficial, back-channel communications that enabled constitutional violations throughout this case:

\begin{itemize}
\item E'Lise's direct recruitment of Officer Ethan Smith (bypassing formal reporting)
\item Smith's off-record harassment of my parents (no documented reports)
\item Detective Brueckner's coordination with YWCA board (undisclosed conflict)
\item Interstate officer coordination without formal warrants or extradition process
\item Suppression of exculpatory evidence through informal institutional channels
\end{itemize}

\item These back-channel communications demonstrate consciousness of wrongdoing---participants knew their actions would not withstand formal scrutiny, so they deliberately operated outside documented, reviewable processes.

\end{enumerate}

\section*{XIX. E'LISE CHARD'S RECRUITMENT OF TYLEEN ROOT: BUSINESS SABOTAGE THROUGH YWCA CLIENTS}

\begin{enumerate}

\item E'Lise Chard expanded her harassment campaign by recruiting Tyleen Root, a client of YWCA Missoula, to engage in a coordinated campaign of slander, defamation, and business sabotage against me. This recruitment of vulnerable domestic violence survivors to serve as proxies for harassment represents an egregious abuse of E'Lise's position and a violation of YWCA's ethical obligations to protect---not exploit---clients.

\textbf{A. Tyleen Root Background and YWCA Connection}

\item Tyleen Root was a client receiving services at YWCA Missoula during the time E'Lise Chard worked there as Operations Manager.

\item I have never met Tyleen Root, have never had any contact with her, and had no knowledge of her existence prior to her harassment campaign.

\item E'Lise recruited Tyleen to serve as a proxy harasser, providing her with false information about me and directing her to engage in a systematic campaign of business sabotage.

\textbf{B. Escalating Pattern of Harassment}

\item Tyleen's harassment followed a deliberate escalation pattern clearly designed to maximize damage while creating plausible deniability for E'Lise:

\item \textbf{Phase 1: Private Direct Messages (Initial Contact)}

Tyleen began by sending private direct messages to me via social media, making false accusations that I had stalked her and her family while she was sheltered at the YWCA.

\begin{itemize}
\item These claims were completely fabricated
\item I have never stalked anyone, particularly not a person I have never met
\item I never responded to any of Tyleen's messages
\item My non-response was a deliberate choice to avoid engagement and escalation with a stranger clearly attempting to antagonize me
\end{itemize}

\item \textbf{Phase 2: Public Posting on Business Social Media}

After failing to provoke a response through private messages, Tyleen escalated to posting public accusations on my business's social media page:

\begin{itemize}
\item Posted on the Facebook page of Pacific Northwest Rural Broadband Alliance (my nonprofit organization)
\item Public accusation: ``The guy who owns this is a stalker and a woman beater''
\item This was posted on a professional business page serving rural communities
\item Designed to destroy my professional reputation and sabotage business relationships
\item Posted with knowledge that it would be seen by clients, partners, and community members
\end{itemize}

\item \textbf{Phase 3: Direct Contact with Page Followers}

When public posting did not provoke the desired response, Tyleen further escalated by directly messaging anyone who interacted with my business page:

\begin{itemize}
\item Sent direct messages to individuals who ``liked,'' commented on, or followed the business page
\item Repeated the false claim that I stalked her and her family while she was at the YWCA
\item Designed to systematically poison all business relationships and community connections
\item Created an atmosphere where any potential client or partner was immediately warned against working with me
\end{itemize}

\item Throughout all three phases of escalation, I never responded to Tyleen Root or attempted any contact with her whatsoever.

\textbf{C. Law Enforcement Response: Protection for Harassers, None for Victims}

\item When I attempted to report Tyleen's harassment and business sabotage to Missoula Police Department, I was subjected to intimidation and threats rather than receiving assistance:

\item \textbf{The Missoula PD Hallway Incident}

\begin{itemize}
\item I went to MPD headquarters to file a report about Tyleen's defamatory social media posts
\item Two officers cornered me in a hallway
\item Officers first ridiculed and laughed at my complaint
\item Officers then yelled at me, escalating to aggressive confrontation
\item When I attempted to reach for my phone to show them screenshots of the harassing messages, one officer yelled: ``DON'T PUT YOUR HANDS IN YOUR POCKET!''
\item The same officer unclipped the latch on his sidearm and placed his hand on his gun, ready to draw
\item I was in imminent fear of being shot by these officers for attempting to report a crime
\item Officers ultimately refused to file a report after this intimidation
\end{itemize}

\item This response demonstrates MPD's institutional alignment with YWCA and continued protection of E'Lise Chard's harassment proxies.

\item \textbf{Dillon Police Department's Professional Response}

\item After determining that Tyleen Root resided in Dillon, Montana, I contacted Dillon Police Department directly.

\item Officer Ian Turnes of Dillon PD immediately took my complaint seriously and recognized Tyleen as someone with a pattern of harassing behavior.

\item Officer Turnes went above and beyond his duties:

\begin{itemize}
\item Immediately recognized Tyleen due to a repeated history of similar harassment, involvement with law enforcement, and Child Protective Services
\item Investigated Tyleen's location (she was actively hiding from Child Protective Services)
\item Located and personally contacted Tyleen
\item Ensured she removed the defamatory posts from my business page
\item Documented the harassment for potential criminal charges
\end{itemize}

\item The stark contrast between Dillon PD's professional response and Missoula PD's threatening intimidation demonstrates that Missoula PD's conduct was not standard law enforcement practice, but rather institutional protection of YWCA-affiliated harassers.

\textbf{D. E'Lise's Abuse of Position and YWCA Liability}

\item E'Lise Chard's recruitment of Tyleen Root represents multiple violations:

\item \textbf{Abuse of Professional Position}

\begin{itemize}
\item E'Lise used her access to YWCA clients to recruit harassment proxies
\item She provided confidential information about me to a client who had no legitimate need for such information
\item She weaponized a domestic violence perpetrator, later identified to be the woman involved with the Arthur Brown case (see: Supplemental Submission, Victim 1: Arthur Brown) exploiting her predisposition for personal vendettas
\item She violated the trust relationship between YWCA and clients receiving services
\end{itemize}

\item \textbf{YWCA Institutional Liability}

\begin{itemize}
\item YWCA employees using client relationships to facilitate harassment
\item Organizational resources and access exploited for personal retaliation
\item No investigation or accountability despite documented abuse
\item Pattern of conduct: E'Lise recruited multiple individuals (Rebecca Pettit, Steve Leutcheman, Tyleen Root) to serve as harassment proxies
\end{itemize}

\item \textbf{Ethical Violations}

\begin{itemize}
\item Violation of professional boundaries between service providers and clients
\item Exploitation of vulnerable populations for personal gain
\item Breach of duty to protect clients from being manipulated into criminal conduct
\item Use of nonprofit resources for private, malicious purposes
\end{itemize}

\textbf{E. Business Damages from Tyleen Root Campaign}

\item Tyleen's harassment campaign, directed by E'Lise Chard, caused quantifiable business damages:

\begin{itemize}
\item Destruction of professional reputation in rural Montana communities
\item Loss of potential contracts and partnerships with individuals who saw the false accusations
\item Damage to Pacific Northwest Rural Broadband Alliance's mission to serve rural communities
\item Psychological trauma from being unable to operate a business without coordinated sabotage
\item Chilling effect on my ability to engage in professional networking or community service
\end{itemize}

\textbf{F. Pattern Recognition: Recruitment of Harassment Proxies}

\item E'Lise's recruitment of Tyleen Root fits an established pattern of enlisting third parties to conduct harassment on her behalf:

\begin{table}[h]
\begin{tabular}{|l|l|l|}
\hline
\textbf{Proxy Recruited} & \textbf{Relationship to E'Lise} & \textbf{Actions Taken} \\
\hline
Rebecca Pettit & YWCA subordinate & False police reports, institutional coordination \\
\hline
Steve Leutcheman & Boyfriend (suspected) & Death threats, repeated harassment \\
\hline
Officer Ethan Smith & YWCA liaison & Year-long surveillance and harassment \\
\hline
Tyleen Root & YWCA client & Business sabotage, social media defamation \\
\hline
Detective Brueckner & YWCA board member & Fraudulent warrant, prosecution coordination \\
\hline
\end{tabular}
\caption{E'Lise Chard's Network of Recruited Harassment Proxies}
\end{table}

\item This pattern demonstrates E'Lise's systematic abuse of every relationship and institutional connection available to her, transforming professional colleagues, romantic partners, law enforcement liaisons, nonprofit clients, and board members into weapons against me.

\textbf{G. Criminal Liability}

\item E'Lise's recruitment of Tyleen Root to engage in defamation and business sabotage constitutes:

\begin{itemize}
\item \textbf{Criminal Solicitation} (MCA § 45-4-102): Commanding or soliciting another to commit a crime (defamation, stalking)
\item \textbf{Criminal Defamation} (MCA § 45-8-212): Knowingly communicating false statements designed to damage reputation
\item \textbf{Conspiracy} (MCA § 45-4-102): Agreement with Tyleen to engage in criminal conduct
\item \textbf{Abuse of Position}: Using professional access to vulnerable populations for criminal purposes
\end{itemize}

\end{enumerate}

\section*{XX. RICO PREDICATES: ORGANIZED CRIMINAL ENTERPRISE}

\begin{enumerate}

\item The coordinated campaign of harassment, false reporting, malicious prosecution, and civil rights violations documented in this declaration constitutes a pattern of racketeering activity under the Racketeer Influenced and Corrupt Organizations Act (RICO), 18 U.S.C. §§ 1961--1968.

\item The enterprise involved in this RICO conspiracy includes:

\begin{itemize}
\item \textbf{Private Citizens}: E'Lise Chard, Danielle Chard, Rebecca Pettit, Steve Leutcheman, Tyleen Root
\item \textbf{Private Nonprofit Organization}: YWCA Missoula (institution and employees)
\item \textbf{Law Enforcement Officers}: Officer Ethan Smith (Missoula PD), Detective Connie Brueckner (Missoula PD), Edmonds Police officers
\item \textbf{Law Enforcement Agencies}: Missoula Police Department, Edmonds Police Department, Seattle Police Department
\item \textbf{Prosecutors}: Missoula County Attorney's Office, King County Prosecutor's Office (multiple jurisdictions)
\end{itemize}

\textbf{A. RICO Enterprise Elements}

\item Under 18 U.S.C. § 1961(4), an ``enterprise'' includes ``any union or group of individuals associated in fact although not a legal entity.'' The coordinated actors in this case constitute such an enterprise through:

\begin{itemize}
\item Common purpose: Suppression of my protected speech and punishment for YWCA complaint
\item Organizational structure: Systematic coordination across institutional boundaries and state lines
\item Longevity: Continuous operation from 2015 through present (10+ years)
\item Interstate scope: Coordinated activities across Washington and Montana
\end{itemize}

\textbf{B. Pattern of Racketeering Activity}

\item Under 18 U.S.C. § 1961(5), a ``pattern of racketeering activity'' requires at least two predicate acts within ten years. The enterprise in this case committed numerous predicate acts that constitute federal crimes under RICO.

\textbf{CHRONOLOGICAL LIST OF RICO PREDICATE ACTS:}

\item \textbf{November 22, 2015 --- Mail and Wire Fraud}

\begin{itemize}
\item \textbf{Predicate}: 18 U.S.C. §§ 1341 (mail fraud), 1343 (wire fraud)
\item \textbf{Act}: Danielle Chard's false report to Edmonds Police claiming I was armed and intended murder-suicide
\item \textbf{Interstate Element}: Communications via telephone and electronic systems across jurisdictions
\item \textbf{Fraudulent Scheme}: False police report designed to provoke lethal law enforcement response
\item \textbf{Harm}: SWAT-style raid placing me at extreme risk of being shot by police
\end{itemize}

\item \textbf{March 2016 --- Extortion Under Color of Official Right}

\begin{itemize}
\item \textbf{Predicate}: 18 U.S.C. § 1951 (Hobbs Act extortion)
\item \textbf{Act}: Coerced guilty plea obtained through threats of extended incarceration and additional charges
\item \textbf{Participants}: Edmonds prosecutors, Edmonds Police, attorney Patricia Fulton (ineffective assistance)
\item \textbf{Harm}: Loss of liberty, criminal record, basis for future harassment
\end{itemize}

\item \textbf{December 2015 -- August 2016 --- Deprivation of Civil Rights Under Color of Law}

\begin{itemize}
\item \textbf{Predicate}: 18 U.S.C. § 242 (deprivation of rights under color of law)
\item \textbf{Acts}: Seattle and Edmonds Police advancement of charges despite documented lack of evidence
\item \textbf{Constitutional Violations}: Fourth Amendment (false arrest), Fifth Amendment (coerced plea), Fourteenth Amendment (due process)
\item \textbf{Result}: All Seattle charges dismissed with prejudice; Edmonds plea obtained through coercion
\end{itemize}

\item \textbf{June 2017 --- Conspiracy Against Rights}

\begin{itemize}
\item \textbf{Predicate}: 18 U.S.C. § 241 (conspiracy against rights)
\item \textbf{Participants}: E'Lise Chard, Officer Ethan Smith, Rebecca Pettit
\item \textbf{Act}: E'Lise's direct recruitment of Officer Smith to surveil and harass me and my parents
\item \textbf{Conspiracy Elements}: Agreement to deprive constitutional rights; overt acts (surveillance, harassment of parents)
\item \textbf{Interstate Element}: Coordination between Washington-based accusers and Montana law enforcement
\end{itemize}

\item \textbf{June 2017 -- May 2018 --- Obstruction of Justice}

\begin{itemize}
\item \textbf{Predicate}: 18 U.S.C. § 1503 (obstruction of justice)
\item \textbf{Act}: Officer Smith's year-long off-record harassment without documented reports
\item \textbf{Purpose}: Prevent oversight and accountability for unconstitutional conduct
\item \textbf{Harm}: Harassment of my elderly parents, contributing to mother's mental health crisis requiring emergency psychiatric hospitalization
\end{itemize}

\item \textbf{June 3, 2018 --- RICO Conspiracy Coordination}

\begin{itemize}
\item \textbf{Date}: My protected First Amendment complaint to YWCA
\item \textbf{Enterprise Response}: Immediate coordination among multiple participants to retaliate
\end{itemize}

\item \textbf{June 12, 2018 --- Mail and Wire Fraud (Perjury in Legal Documents)}

\begin{itemize}
\item \textbf{Predicate}: 18 U.S.C. §§ 1341, 1343
\item \textbf{Act}: E'Lise Chard's protection order petition containing demonstrably false statements
\item \textbf{False Claims}: Extensive criminal drug history (I have none); threats to children; ongoing contact (admitted under oath only 3 emails)
\item \textbf{Interstate Element}: Filed in Montana based on Washington coordination
\item \textbf{Purpose}: Create basis for criminal charges in retaliation for YWCA complaint
\end{itemize}

\item \textbf{June 2018 --- Perjury}

\begin{itemize}
\item \textbf{Predicate}: 18 U.S.C. § 1621 (perjury)
\item \textbf{Act}: E'Lise's contradictory sworn testimony before Judge Warren
\item \textbf{Contradiction}: Written petition alleged extensive ongoing harassment; sworn testimony admitted only 3 emails
\item \textbf{Materiality}: False statements formed entire basis for criminal prosecution
\end{itemize}

\item \textbf{June -- September 2018 --- Conspiracy Against Rights (YWCA-Police Coordination)}

\begin{itemize}
\item \textbf{Predicate}: 18 U.S.C. § 241
\item \textbf{Participants}: Detective Connie Brueckner (YWCA board member), E'Lise Chard, YWCA institutional leadership
\item \textbf{Conspiracy}: Undisclosed conflict of interest enabling retaliatory prosecution
\item \textbf{Overt Acts}: Fraudulent search warrant, home raid, seizure of work equipment, GPS monitoring
\item \textbf{Constitutional Violations}: First Amendment (retaliation), Fourth Amendment (fraudulent warrant), Fourteenth Amendment (due process)
\end{itemize}

\item \textbf{July 2018 --- Deprivation of Civil Rights (Home Raid)}

\begin{itemize}
\item \textbf{Predicate}: 18 U.S.C. § 242
\item \textbf{Act}: Warrantless home entry during business meeting; seizure of all electronics
\item \textbf{Basis}: Search warrant obtained through Detective Brueckner's undisclosed YWCA conflict and Officer Smith's fabricated statements
\item \textbf{Constitutional Violations}: Fourth Amendment (fraudulent warrant basis), Brady (suppression of exculpatory evidence)
\item \textbf{Harm}: Loss of \$150,000 annual IT contract; career destruction
\end{itemize}

\item \textbf{August 14, 2018 --- Obstruction of Justice (Suppression of Exculpatory Evidence)}

\begin{itemize}
\item \textbf{Predicate}: 18 U.S.C. § 1503
\item \textbf{Act}: YWCA's deliberate suppression of Witness PS character reference
\item \textbf{Participants}: YWCA institutional leadership, Detective Brueckner (in her capacity as board member)
\item \textbf{Exculpatory Evidence Suppressed}: Firsthand testimony contradicting E'Lise's claims; documentation of my non-violent character
\item \textbf{Brady Violation}: Concealment of material evidence favorable to defense
\end{itemize}

\item \textbf{August 2018 --- Interstate Communication of False Information}

\begin{itemize}
\item \textbf{Predicate}: 18 U.S.C. §§ 1341, 1343 (mail and wire fraud)
\item \textbf{Act}: E'Lise communicated my custody status to Danielle in Washington
\item \textbf{Purpose}: Coordinate warrantless arrest through interstate back-channel
\item \textbf{Result}: Unlawful arrest, days ``lost in jail'' without charges
\end{itemize}

\item \textbf{August 2018 --- Deprivation of Rights (Warrantless Arrest)}

\begin{itemize}
\item \textbf{Predicate}: 18 U.S.C. § 242
\item \textbf{Act}: Montana officer's warrantless home entry and arrest based on fabricated extradition claim
\item \textbf{Coordination}: Edmonds officer's unofficial call to Montana officer at Danielle's request
\item \textbf{Constitutional Violations}: Fourth Amendment (warrantless arrest without probable cause), Fourteenth Amendment (detention without charges)
\item \textbf{Harm}: Multiple days detained without lawful basis
\end{itemize}

\item \textbf{September 2018 --- Mail and Wire Fraud (Charging Documents)}

\begin{itemize}
\item \textbf{Predicate}: 18 U.S.C. §§ 1341, 1343
\item \textbf{Act}: Stalking charging documents based entirely on Officer Smith's fabricated statements
\item \textbf{Fabrications}: Non-existent phone conversations with parents; false mental health claims
\item \textbf{Interstate Element}: Charging documents relied on Washington-Montana coordination
\item \textbf{Result}: 18 months of prosecution; all charges ultimately dismissed with prejudice
\end{itemize}

\item \textbf{September 2018 -- 2019 --- Interstate Coordination for Harassment}

\begin{itemize}
\item \textbf{Predicate}: 18 U.S.C. § 241 (conspiracy against rights)
\item \textbf{Participants}: Washington prosecutors (Seattle, Lynnwood, Edmonds Municipal, Edmonds Superior), Montana prosecutors, law enforcement across jurisdictions
\item \textbf{Coordination}: Simultaneous charges in seven different courts across two states
\item \textbf{Purpose}: Overwhelm defense capacity; drain financial resources; maximize pretrial punishment
\item \textbf{Interstate Commerce Impact}: Interference with multi-state IT consulting business
\end{itemize}

\item \textbf{February 2020 --- Witness Tampering}

\begin{itemize}
\item \textbf{Predicate}: 18 U.S.C. § 1512 (tampering with witness, victim, or informant)
\item \textbf{Act}: Edmonds Police officer threatened romantic partner with loss of child custody unless she testified against me
\item \textbf{Threat}: ``You will lose custody of your child and never see your son again''
\item \textbf{Purpose}: Coerce testimony; force guilty plea
\item \textbf{Result}: Coerced guilty pleas in September 2019 to protect partner and child
\item \textbf{Constitutional Violations}: Fifth Amendment (coerced self-incrimination), Sixth Amendment (denial of fair trial)
\end{itemize}

\item \textbf{2020 -- 2024 --- Continuing Conspiracy and Business Sabotage}

\begin{itemize}
\item \textbf{Predicate}: 18 U.S.C. §§ 241, 1951 (conspiracy; interference with commerce by threats or violence)
\item \textbf{Acts}: E'Lise's recruitment of Tyleen Root (YWCA client) to conduct business sabotage campaign
\item \textbf{Business Impact}: Social media defamation; direct contact with clients and partners; loss of \$60,000 annual St. Regis School District contract
\item \textbf{Interstate Commerce}: Interference with multi-state nonprofit broadband development
\end{itemize}

\item \textbf{2018 -- 2025 --- Pattern of Obstruction (Missoula PD Institutional)}

\begin{itemize}
\item \textbf{Predicate}: 18 U.S.C. § 1503 (obstruction of justice)
\item \textbf{Acts}: Repeated refusal to investigate harassment complaints; intimidation when attempting to file reports; threats against me for seeking accountability
\item \textbf{Institutional Pattern}: Protecting YWCA-affiliated harassers while denying victim protection
\item \textbf{Example}: Hallway confrontation with officers when reporting Tyleen Root harassment
\end{itemize}

\textbf{C. RICO Damages}

\item Under 18 U.S.C. § 1964(c), any person injured in business or property by reason of a RICO violation is entitled to treble damages (three times actual damages) plus attorney fees and costs.

\item \textbf{Actual Damages Caused by RICO Enterprise:}

\begin{table}[h]
\begin{tabular}{|l|r|}
\hline
\textbf{Damage Category} & \textbf{Amount} \\
\hline
Lost IT Career (2018--2025) & \$2,100,000 \\
Contract Losses (specific contracts) & \$680,000 \\
Legal Fees (multiple jurisdictions) & \$387,000 \\
Medical/Psychological Treatment & \$148,000 \\
Family Medical Expenses (parents) & \$76,000 \\
Business Development Losses & \$30,000 \\
\hline
\textbf{Total Actual Damages} & \$3,421,000 \\
\hline
\textbf{Treble Damages (RICO)} & \$10,263,000 \\
\hline
\end{tabular}
\caption{RICO Damages Calculation}
\end{table}

\item These damages do not include:

\begin{itemize}
\item Punitive damages (available under state law claims)
\item Non-economic damages (pain and suffering, emotional distress)
\item Future economic losses (diminished earning capacity)
\item Attorney fees and costs (recoverable under RICO)
\end{itemize}

\textbf{D. Federal Jurisdiction and Interstate Commerce}

\item RICO jurisdiction is established through multiple interstate elements:

\begin{itemize}
\item Coordination between Washington and Montana actors
\item Use of mail and wire communications across state lines
\item Impact on interstate commerce (IT consulting business, nonprofit broadband development)
\item Multi-state legal proceedings (7 courts, 2 states)
\item Interstate travel restrictions imposed through GPS monitoring
\end{itemize}

\textbf{E. Conclusion: RICO Enterprise Liability}

\item The coordinated campaign documented in this declaration constitutes a classic RICO enterprise:

\begin{itemize}
\item \textbf{Longevity}: 10+ years of continuous operation (2015--2025)
\item \textbf{Coordination}: Systematic collaboration across institutional and jurisdictional boundaries
\item \textbf{Multiple Predicate Acts}: At least 20 separate federal crimes committed in furtherance of enterprise
\item \textbf{Common Purpose}: Punishment of protected speech and suppression of institutional accountability
\item \textbf{Interstate Scope}: Coordinated activities across state lines
\item \textbf{Quantifiable Damages}: \$3.4+ million in direct losses; \$10+ million in treble damages
\end{itemize}

\item This RICO enterprise represents the weaponization of law enforcement, prosecutors, and private nonprofit organizations to achieve personal vendettas and institutional protection at the expense of constitutional rights and the rule of law.

\end{enumerate}
\section*{XXI. SUMMARY AND REQUEST FOR ACTION}

\begin{enumerate}

\item Taken together, the events described in this declaration show a multi-year, multi-jurisdictional pattern in which E'Lise and Danielle Chard repeatedly made false statements, committed perjury, and engaged in criminal defamation, with assistance and cover from certain law enforcement officers and institutional actors.

\item I respectfully request that state and federal authorities open or expand criminal investigations into:

\begin{itemize}
  \item Perjury and criminal defamation by E'Lise Chard;
  \item Perjury and false reporting by Danielle Chard;
  \item Witness intimidation and coercion by Edmonds Police officers and other officials;
  \item Suppression of exculpatory evidence and conflicts of interest by YWCA Missoula and Detective Connie Brueckner;
  \item Conspiracy to violate my civil rights under 18 U.S.C. \S 241 and related Montana statutes.
\end{itemize}

\item I am willing to fully cooperate with any such investigation, provide documentation, identify additional witnesses, and testify under oath before a grand jury or at trial.

\end{enumerate}

\section*{XXII. DECLARATION AND SIGNATURE}

I declare under penalty of perjury under the laws of the State of Montana and the United States of America that the foregoing is true and correct to the best of my knowledge, information, and belief. I understand that providing false information in this declaration may subject me to criminal prosecution for perjury.

I acknowledge that I accepted guilty pleas in 2016 and 2019 under coercive circumstances as detailed herein. These pleas do not reflect my actual guilt but rather the unconstitutional pressures applied by law enforcement and prosecutors, including threats directed at an innocent third party and her child. I am prepared to testify under oath regarding the involuntary nature of these pleas and the governmental misconduct that produced them.

I am available for interview, deposition, grand jury testimony, and trial testimony regarding all matters in this declaration.

\vspace{1cm}

\noindent Executed this \rule{1in}{0.5pt} day of \rule{1.5in}{0.5pt}, 2025, at \rule{2in}{0.5pt}, Montana.

\vspace{1cm}

\noindent\rule{3.5in}{0.5pt}\\
\textbf{Elvis Ryland Nuno}\\
Declarant

\vspace{1cm}

\noindent \textbf{NOTARY ACKNOWLEDGMENT}

\vspace{0.5cm}

\noindent SUBSCRIBED AND SWORN TO (or affirmed) before me on this \rule{1in}{0.5pt} day of \rule{1.5in}{0.5pt}, 2025, by Elvis Ryland Nuno, who is personally known to me or has provided satisfactory evidence of identity.

\vspace{1cm}

\noindent\rule{3.5in}{0.5pt}\\
Notary Public for the State of Montana\\
My Commission Expires: \rule{1.5in}{0.5pt}

\end{document}
